% !TEX root = main.tex
\section*{2. 어떤 설명을 추가하거나 강조하면 좋은가?}
\enlargethispage{2\baselineskip}
이 부분부터는 원문 요약이 아니라, 지금까지의 대화를 바탕으로 한 변경 제안입니다.

\subsection*{제품에서 과거로 찾는 경우도 설명하기}
현재 1.2는 원료의 증빙에 문제가 발견된 상황에서 시작합니다.
그래서 후속 기록을 찾아 동결하는 이유는 보이지만, 완제품에서 원료·공정을 거슬러 올라가는 이유는 약하게 드러납니다.

\textbf{완제품에 문제가 발견됐지만 어느 원료나 공정에서 비롯됐는지 아직 모르는 상황}도 함께 제시하면 좋겠습니다.
Backward Tracing은 그 제품에 사용된 입력과 공정을 조사할 이력을 제공합니다.
원인을 자동 판정하는 기능이라고 설명할 필요는 없습니다.

반대로 문제가 있는 원료를 이미 알고 있다면, Forward Tracing으로 그 원료를 사용한 후속 기록을 찾습니다.
이렇게 두 상황을 설명한 뒤, 사용 중단이 승인되면 아직 소비되지 않은 대상을 Freeze한다고 이어갈 수 있습니다.

\subsection*{Forward Tracing이 왜 더 까다로운지 설명하기}
거래를 만든 사람은 자신이 어떤 입력을 사용했는지 알고 있습니다.
따라서 Backward Tracing에 필요한 부모 정보를 암호화해 남기는 구조를 생각할 수 있습니다.

Forward Tracing은 생성된 객체가 나중에 어느 거래에서 소비됐는지 연결해야 합니다.
현재 Note 구조에서 그 연결에 쓰는 nf는 소유자의 비밀값으로 계산합니다.

Sender가 Receiver의 Note를 직접 만든다면, Sender는 Receiver의 비밀값을 모르므로 그 Note의 미래 nf를 계산할 수 없습니다.
반대로 Sender도 nf를 계산할 수 있게 바꾸면, Sender가 공개 원장에서 Receiver의 다음 소비를 알아낼 수 있습니다.

\textbf{보낸 사람에게 이후 소비를 드러내지 않으면서, 승인된 Auditor는 추적을 이어가게 만드는 것이 설명해야 할 설계 문제입니다.}
이것이 다른 방법으로는 불가능하다는 주장은 아니며, 현재 소유권·nf 관계에서 발생하는 어려움입니다.

\subsection*{Voucher를 왜 사용하는지 설명하기}
Voucher는 단순히 전달 절차를 한 번 늘리는 장치가 아닙니다.
Receiver가 Proceed할 때 자기 Note를 만들고, 자신만 계산할 수 있는 미래 nf를 Auditor용으로 암호화해 남기도록 합니다.

Sender는 Voucher의 해결 여부와 Proceed에서 생성된 cm까지는 연결할 수 있습니다.
숨기려는 것은 \textbf{그 Receiver Note가 이후 어느 거래에서 소비되는지}입니다.

또한 Receiver가 인수하기 전에는 Sender가 정해진 조건 안에서 Recall할 수 있습니다.
Introduction에서는 이 두 목적을 설명하고, 상세 수식과 deadline 조건은 본문으로 넘기면 됩니다.

\subsection*{Contribution은 기능과 확인한 결과를 구분해서 설명하기}
첫 번째 후보는 비공개 State를 가진 객체로 이전·분할·병합·가공을 표현하고, 그 결과를 Claim 발급에 연결하는 프로토콜입니다.
이는 자산 한 번의 전달보다 넓은 공급망 동작을 다룬다는 설명입니다.

두 번째 후보는 실제 거래와 일치하는 감사 정보, Voucher를 이용한 이전, 승인된 양방향 추적과 소비 차단을 함께 구성하는 것입니다.
개별 기능이 이미 존재한다는 점을 인정하면서, 무엇을 함께 만족하도록 설계했는지 밝혀야 합니다.

세 번째는 그 설계에 대해 실제로 확인한 결과입니다.
현재 초안의 보안 분석·기호 모델 검사와, 우리가 수행한 구현·성능 실험을 같은 근거처럼 섞어서는 안 됩니다.

\textbf{실험을 논문에 포함하려면 먼저 논문의 키 구조·Event 규칙과 실험 구현이 일치하는지 확인해야 합니다.}
위 항목이 선행연구에 비해 어떤 새 기여인지는 별도의 비교가 필요합니다.
